\documentclass[12pt,a4paper,oneside]{book}

% =====================================================================
% PACOTES ESSENCIAIS
% =====================================================================
\usepackage[brazilian]{babel}
\usepackage[T1]{fontenc}
\usepackage[utf8]{inputenc}
\usepackage{amsmath}
\usepackage{amssymb}
\usepackage{amsthm}
\usepackage{geometry}
\usepackage{fancyhdr}
\usepackage{hyperref}
\usepackage{booktabs}
\usepackage{mathtools}

% =====================================================================
% GEOMETRIA
% =====================================================================
\geometry{
    left=2.5cm,
    right=2.5cm,
    top=3cm,
    bottom=3cm,
    headheight=14.5pt
}

% =====================================================================
% CABECALHO E RODAPE
% =====================================================================
\pagestyle{fancy}
\fancyhf{}
\fancyhead[L]{\textit{Capitulo 3: Escala, Escalabilidade e Fractais}}
\fancyhead[R]{\thepage}
\fancyfoot[C]{\small Principios de Modelagem Matematica}
\renewcommand{\headrulewidth}{0.4pt}
\renewcommand{\footrulewidth}{0.4pt}

% =====================================================================
% AMBIENTES
% =====================================================================
\theoremstyle{definition}
\newtheorem{definition}{Definicao}[chapter]
\newtheorem{example}[definition]{Exemplo}
\newtheorem{remark}[definition]{Observacao}
\newtheorem{problem}{Problema}

% =====================================================================
% DOCUMENTO
% =====================================================================
\title{\textbf{Capitulo 3: Escala, Escalabilidade e Fractais}\\[0.4cm]
\large Principios de Modelagem Matematica}
\author{Baseado em Dym, C. L.}
\date{2026-04-14}

\begin{document}

\maketitle
\tableofcontents
\newpage

\chapter{Escala, Escalabilidade e Fractais}

\section{Ideia central do capitulo}

Neste capitulo, o foco e entender como um sistema muda quando alteramos seu tamanho caracteristico.
Esse tipo de pergunta aparece em engenharia, fisica, biologia e economia:

\begin{itemize}
    \item Se dobramos o tamanho, o efeito tambem dobra?
    \item Quais grandezas crescem mais rapido e se tornam dominantes?
    \item Existe uma lei de potencia do tipo $Y \propto X^\alpha$?
\end{itemize}

A analise de escala permite identificar termos relevantes sem resolver o problema completo.

\section{Escala e analise por ordem de grandeza}

\begin{definition}[Escala caracteristica]
Escala caracteristica e o valor tipico de uma variavel no problema, como comprimento $L$, velocidade $U$, tempo $T$ ou concentracao $C$.
\end{definition}

\begin{definition}[Ordem de grandeza]
Dizemos que duas quantidades sao da mesma ordem quando diferem por um fator pequeno (tipicamente menor que 10).
\end{definition}

\subsection{Nao dimensionalizacao rapida}

Considere a equacao de transporte 1D:
\begin{equation}
\frac{\partial c}{\partial t} + u\frac{\partial c}{\partial x} = D\frac{\partial^2 c}{\partial x^2}.
\end{equation}

Definimos variaveis adimensionais:
\begin{equation}
\hat{x} = \frac{x}{L}, \quad \hat{t} = \frac{t}{T}, \quad \hat{c} = \frac{c}{C_0}, \quad \hat{u} = \frac{u}{U}.
\end{equation}

Escolhendo $T=L/U$, obtemos:
\begin{equation}
\frac{\partial \hat{c}}{\partial \hat{t}} + \hat{u}\frac{\partial \hat{c}}{\partial \hat{x}} = \frac{1}{Pe}\frac{\partial^2 \hat{c}}{\partial \hat{x}^2},
\end{equation}
com numero de Peclet:
\begin{equation}
Pe = \frac{UL}{D}.
\end{equation}

\begin{remark}
Se $Pe \gg 1$, adveccao domina; se $Pe \ll 1$, difusao domina.
\end{remark}

\section{Leis de escala (power laws)}

Em muitos sistemas, observa-se:
\begin{equation}
Y = kX^\alpha,
\end{equation}
onde $\alpha$ e o expoente de escala.

Tomando logaritmo:
\begin{equation}
\log Y = \log k + \alpha \log X,
\end{equation}
logo, em grafico log-log, a inclinacao e $\alpha$.

\subsection{Exemplo 1: area e volume}

Para um corpo com comprimento caracteristico $L$:
\begin{equation}
A \propto L^2, \qquad V \propto L^3.
\end{equation}

Se aumentamos $L$ por fator 3:
\begin{equation}
A' = 3^2A = 9A, \qquad V' = 3^3V = 27V.
\end{equation}

\textbf{Conclusao}: volume cresce mais rapido que area. Isso explica, por exemplo, limites de troca termica em objetos grandes.

\subsection{Exemplo 2: periodo de pendulo}

Para pequenas amplitudes:
\begin{equation}
T = 2\pi\sqrt{\frac{\ell}{g}}.
\end{equation}

Logo:
\begin{equation}
T \propto \ell^{1/2}.
\end{equation}

Se $\ell$ quadruplica, o periodo dobra.

\section{Escalabilidade de modelos}

Escalabilidade significa como desempenho, custo ou erro variam com o tamanho do sistema.

\subsection{Exemplo 3: custo computacional}

Suponha um algoritmo com custo:
\begin{equation}
C(N) = aN^2.
\end{equation}

Se o tamanho de entrada dobra ($N \to 2N$):
\begin{equation}
C(2N) = a(2N)^2 = 4aN^2.
\end{equation}

Ou seja, dobrar o problema quadruplica o custo.

\subsection{Exemplo 4: transferencia de calor em objeto geometricamente similar}

Potencia termica aproximada por conveccao:
\begin{equation}
\dot{Q} \sim hA\Delta T.
\end{equation}

Como $A \propto L^2$, entao:
\begin{equation}
\dot{Q} \propto L^2.
\end{equation}

Energia armazenada (aprox. proporcional ao volume):
\begin{equation}
E \propto L^3.
\end{equation}

Tempo caracteristico de resfriamento:
\begin{equation}
\tau \sim \frac{E}{\dot{Q}} \propto L.
\end{equation}

Objetos maiores tendem a resfriar mais lentamente.

\section{Escala isometrica e alometrica}

Quando todos os comprimentos de um objeto crescem proporcionalmente, chamamos de
\textbf{escala isometrica}. Quando a proporcao entre dimensoes muda, chamamos de
\textbf{escala alometrica}.

\subsection{Escala isometrica: lei do quadrado-cubo}

Na escala isometrica, se o comprimento caracteristico $L$ e multiplicado por fator $\lambda$:
\begin{equation}
L' = \lambda L, \qquad A' = \lambda^2 A, \qquad V' = \lambda^3 V.
\end{equation}

Consequencia biologica importante: um animal duas vezes maior tem 4 vezes mais superficie
(respiracao, troca termica) mas 8 vezes mais massa. Isso impoe limites biologicos ao tamanho maximo.

\subsection{Lei de Kleiber: escala alometrica do metabolismo}

Em escala alometrica, os expoentes diferem dos isometricos. Um exemplo classico e a
\textbf{Lei de Kleiber} (1932), que relaciona a taxa metabolica basal $q_0$ com a massa
corporal $M$ de mamiferos:
\begin{equation}
q_0 \approx 70\,M^{3/4} \qquad [\text{kcal/dia},\; M \text{ em kg}].
\end{equation}

Se a escala fosse isometrica (proporcional a area superficial), esperariamos $q_0 \propto M^{2/3}$.
O expoente $3/4$ observado e maior, indicando escala alometrica positiva.

\begin{remark}
Um elefante ($M \sim 5000\,\text{kg}$) tem metabolismo total muito maior que um rato
($M \sim 0.02\,\text{kg}$), mas seu metabolismo especifico (por kg) e muito menor.
O rato queima energia mais rapido por unidade de massa.
\end{remark}

\subsection{Tabela de leis alometricas em mamiferos}

\begin{center}
\begin{tabular}{lll}
\toprule
\textbf{Grandeza} & \textbf{Escala isometrica esperada} & \textbf{Escala observada} \\
\midrule
Taxa metabolica $q_0$ & $M^{2/3}$ & $M^{3/4}$ \\
Frequencia cardiaca $f_c$ & $M^{-1/3}$ & $M^{-1/4}$ \\
Tempo de vida $t_v$ & $M^{1/3}$ & $M^{1/4}$ \\
Volume de sangue $V_s$ & $M^{1}$ & $M^{1.0}$ (isometrico) \\
\bottomrule
\end{tabular}
\end{center}

\begin{remark}
Em grafico log-log, a Lei de Kleiber aparece como uma reta com inclinacao $3/4$,
abrangendo desde camundongos ($\sim 10^{-2}$ kg) ate baleias ($\sim 10^5$ kg) ---
uma faixa de 7 ordens de grandeza, confirmando a universalidade da lei de potencia.
\end{remark}

\section{Fractais e dimensao fractal}

Fractais sao estruturas com auto-similaridade: partes se parecem com o todo em diferentes escalas.

\begin{definition}[Dimensao fractal por auto-similaridade]
Se um objeto e composto por $N$ copias reduzidas por fator $r$ em cada escala,
a dimensao fractal e:
\begin{equation}
D = \frac{\log N}{\log(1/r)}.
\end{equation}
\end{definition}

\subsection{Exemplo 5: Curva de Koch}

Na Curva de Koch, cada segmento e substituido por $N=4$ segmentos com reducao $r=1/3$:
\begin{equation}
D = \frac{\log 4}{\log 3} \approx 1.2619.
\end{equation}

A curva tem dimensao maior que uma linha ($D=1$), mas menor que uma superficie ($D=2$).

\subsection{Exemplo 6: Triangulo de Sierpinski}

No Triangulo de Sierpinski: $N=3$ e $r=1/2$:
\begin{equation}
D = \frac{\log 3}{\log 2} \approx 1.585.
\end{equation}

\section{Roteiro pratico para estudar o capitulo 3}

\begin{enumerate}
    \item Identifique as variaveis e escalas caracteristicas do problema.
    \item Nao dimensionalize as equacoes.
    \item Extraia numeros adimensionais (Re, Pe, etc.).
    \item Compare ordens de grandeza para eliminar termos pequenos.
    \item Verifique se ha lei de potencia em grafico log-log.
    \item Quando houver auto-similaridade, estime dimensao fractal.
\end{enumerate}

\section{Exercicios curtos}

\begin{problem}
Se $Y \propto X^{1.7}$, por qual fator $Y$ muda quando $X$ dobra?
\end{problem}

\begin{problem}
Um objeto geometrico similar tem aresta multiplicada por 5. Como mudam area e volume?
\end{problem}

\begin{problem}
Em uma malha numerica 3D, o numero de celulas escala como $N \propto (L/\Delta x)^3$. Se $\Delta x$ cai pela metade, qual o fator de aumento em $N$?
\end{problem}

% =====================================================================
% EXEMPLO TIPO PROVA
% =====================================================================
\section{Exemplo tipo prova: dimensao fractal}

\begin{problem}[Prova -- IIo semestre 2013 e 2025]
Determine a dimensao fractal dos seguintes conjuntos auto-similares:
\begin{enumerate}
    \item \textbf{Curva tipo Koch}: a cada iteracao, cada segmento e substituido por
    4 copias de si mesmo com escala $1/3$.
    \item \textbf{Piramide tipo Sierpinski}: a cada iteracao, a figura e composta por
    4 copias em escala $1/2$.
\end{enumerate}
\end{problem}

\textbf{Solucao}: A formula geral para fractal auto-similar e:
\begin{equation}
D_f = \frac{\ln N_p}{\ln(1/p)},
\end{equation}
onde $N_p$ e o numero de copias por iteracao e $p$ e o fator de escala linear.

\textbf{1. Curva tipo Koch}: $N_p = 4$, $p = 1/3$:
\begin{equation}
D_f = \frac{\ln 4}{\ln 3} = \frac{2\ln 2}{\ln 3} \approx \frac{1{,}386}{1{,}099} \approx \boxed{1{,}26}.
\end{equation}
Como $1 < D_f < 2$: o objeto e mais complexo que uma curva 1D, mas nao preenche o plano 2D.

\textbf{2. Piramide tipo Sierpinski}: $N_p = 4$, $p = 1/2$:
\begin{equation}
D_f = \frac{\ln 4}{\ln 2} = \frac{2\ln 2}{\ln 2} = 2 \implies \boxed{D_f = 2}.
\end{equation}

\textbf{Verificacao iterativa do item 2}:
\begin{center}
\begin{tabular}{ccc}
\toprule
Iteracao $k$ & Numero de copias $N$ & Escala $p^k$ \\
\midrule
1 & 4 & $1/2$ \\
2 & 16 & $1/4$ \\
$n$ & $4^n$ & $1/2^n$ \\
\bottomrule
\end{tabular}
\end{center}
$D_f = \displaystyle\lim_{k\to\infty}\frac{\ln 4^k}{\ln 2^k} = \frac{\ln 4}{\ln 2} = 2.\quad\checkmark$

\begin{remark}
$D_f = 2$ para a piramide Sierpinski e surpreendente: apesar da aparencia rendada
com infinitos buracos, ela tem a mesma dimensao fractal de um plano solido.
Isso ocorre porque o numero de copias ($N=4$) cresce tao rapido quanto a reducao de escala
($p=1/2$ em 2D implica area $\propto p^2 = 1/4$), resultando em ``densidade'' plena.
\end{remark}

\section{Resumo final}

O capitulo 3 mostra que escalar nao e apenas aumentar tamanho: e entender como os mecanismos mudam de importancia.
A linguagem matematica dessa ideia aparece em:
\begin{itemize}
    \item leis de potencia,
    \item numeros adimensionais,
    \item eliminacao de termos por ordem de grandeza,
    \item geometria fractal e dimensao nao inteira.
\end{itemize}

Essas tecnicas tornam a modelagem mais robusta, economica e explicativa.

% =====================================================================
% RESUMO: EXPLICACAO E EXEMPLO DE CADA CONCEITO
% =====================================================================
\section{Resumo: explicacao e exemplo de cada conceito}

% -------------------------------------------------------------------
\subsection*{1. Escala caracteristica e ordem de grandeza}

\textbf{O que e:}
Escala caracteristica e o valor tipico que uma variavel assume no problema
(ex.: comprimento $L$, velocidade $U$, tempo $T$).
Duas grandezas sao da mesma ordem de grandeza quando seu quociente e proximo de 1
(nao difere por mais de uma potencia de 10).

\textbf{Para que serve:}
Permite comparar termos de uma equacao e descartar os pequenos antes de resolver.

\textbf{Exemplo:}
Na equacao do calor
$\rho c_p \dfrac{\partial T}{\partial t} = k\nabla^2 T$,
as escalas tipicas sao $\Delta T$, $L$ e $\tau$.
Estimando cada termo:
\[
\rho c_p \frac{\Delta T}{\tau} \sim k\frac{\Delta T}{L^2}
\quad\Rightarrow\quad
\tau \sim \frac{\rho c_p L^2}{k}.
\]
Sem resolver a EDP, ja sabemos que o tempo caracteristico de difusao termica
escala com $L^2$.

% -------------------------------------------------------------------
\subsection*{2. Nao dimensionalizacao e numero de Peclet}

\textbf{O que e:}
Substituir variaveis fisicas por versoes adimensionais ($\hat{x}=x/L$, etc.)
para revelar os parametros que realmente controlam o fenomeno.

\textbf{Para que serve:}
Reduz o numero de parametros do problema e mostra qual mecanismo domina.

\textbf{Exemplo:}
Para a equacao de transporte de soluto
$\partial_t c + u\,\partial_x c = D\,\partial_{xx} c$,
apos nao dimensionalizar com $\hat{x}=x/L$, $\hat{t}=tU/L$, $\hat{c}=c/C_0$:
\[
\partial_{\hat{t}}\hat{c} + \hat{u}\,\partial_{\hat{x}}\hat{c}
= \frac{1}{Pe}\,\partial_{\hat{x}\hat{x}}\hat{c},
\qquad Pe = \frac{UL}{D}.
\]
Se $Pe = 10^4$ (agua em rio), a difusao e desprezivel e o transporte e puramente advectivo.
Se $Pe = 10^{-3}$ (difusao molecular em gel), a adveccao e desprezivel.

% -------------------------------------------------------------------
\subsection*{3. Lei de potencia (power law)}

\textbf{O que e:}
Relacao do tipo $Y = kX^\alpha$, onde $\alpha$ e o expoente de escala.

\textbf{Para que serve:}
Descreve sistemas que obedecem auto-similaridade ou escalonamento em varias ordens
de grandeza. Em grafico log-log, aparece como reta de inclinacao $\alpha$.

\textbf{Exemplo:}
O periodo de um pendulo simples obedece $T \propto \ell^{1/2}$.
Se medirmos $T$ para varios valores de $\ell$ e plotarmos $\log T \times \log \ell$,
obtemos uma reta com inclinacao $1/2$, confirmando a lei de potencia sem precisar
conhecer $2\pi/\sqrt{g}$ a priori.

\begin{center}
\begin{tabular}{ccc}
\toprule
$\ell$ (m) & $T$ (s) & $\log_{10} T$ \\
\midrule
0.25 & 1.00 & 0.00 \\
1.00 & 2.00 & 0.30 \\
4.00 & 4.00 & 0.60 \\
\bottomrule
\end{tabular}
\end{center}
Inclinacao no grafico log-log: $\Delta(\log T)/\Delta(\log \ell) = 0.30/0.60 = 0.5$. $\checkmark$

% -------------------------------------------------------------------
\subsection*{4. Escalabilidade: lei do quadrado-cubo}

\textbf{O que e:}
Quando um objeto e ampliado por fator $\lambda$ em todas as dimensoes,
area cresce com $\lambda^2$ e volume com $\lambda^3$.

\textbf{Para que serve:}
Explica limites fisicos e biologicos: por que animais gigantes nao existem,
por que microprocessadores aquecem mais ao miniaturizar, etc.

\textbf{Exemplo:}
Um cubo de lado $L=1\,\text{m}$ tem $A=6\,\text{m}^2$ e $V=1\,\text{m}^3$.
Se dobrarmos o lado ($\lambda=2$):
\[
A' = 4\times 6 = 24\,\text{m}^2,\qquad V' = 8\times 1 = 8\,\text{m}^3.
\]
A razao $A/V$ caiu de 6 para 3: o objeto maior troca calor proporcionalmente menos
por unidade de volume, ou seja, resfria mais devagar.

% -------------------------------------------------------------------
\subsection*{5. Escala alometrica e Lei de Kleiber}

\textbf{O que e:}
Escala alometrica ocorre quando os expoentes reais diferem dos esperados pela
geometria isometrica. A Lei de Kleiber e o exemplo canonico:
\[
q_0 \approx 70\,M^{3/4} \qquad [\text{kcal/dia}].
\]

\textbf{Para que serve:}
Revela que biologias, redes vasculares e sistemas complexos nao escalam
simplesmente com area ($M^{2/3}$), mas seguem leis proprias.

\textbf{Exemplo:}
\begin{center}
\begin{tabular}{lcc}
\toprule
Animal & Massa $M$ (kg) & $q_0$ previsto por $70M^{3/4}$ (kcal/dia) \\
\midrule
Rato       & 0{,}02 & $\approx 3{,}5$ \\
Gato       & 3{,}0  & $\approx 152$  \\
Ser humano & 70{,}0 & $\approx 1680$ \\
Elefante   & 4000   & $\approx 44000$\\
\bottomrule
\end{tabular}
\end{center}
O metabolismo especifico (kcal/dia/kg) do rato e $\approx 175$ vezes maior que
o do elefante --- consequencia direta do expoente $3/4 < 1$.

% -------------------------------------------------------------------
\subsection*{6. Dimensao fractal}

\textbf{O que e:}
Medida de complexidade geometrica de objetos auto-similares.
Para um fractal composto por $N$ copias reduzidas por fator $r$:
\[
D = \frac{\log N}{\log(1/r)}.
\]

\textbf{Para que serve:}
Quantifica estruturas que nao tem dimensao inteira (nem curva nem superficie):
costas litoraneas, flocos de neve, redes de rios, pulmoes, etc.

\textbf{Exemplo comparativo:}
\begin{center}
\begin{tabular}{llccc}
\toprule
Fractal & Regra de construcao & $N$ & $r$ & $D$ \\
\midrule
Segmento (referencia) & 2 copias, escala $1/2$ & 2 & $1/2$ & 1{,}00 \\
Curva de Koch         & 4 copias, escala $1/3$ & 4 & $1/3$ & 1{,}26 \\
Triangulo de Sierpinski & 3 copias, escala $1/2$ & 3 & $1/2$ & 1{,}58 \\
Quadrado (referencia) & 4 copias, escala $1/2$ & 4 & $1/2$ & 2{,}00 \\
\bottomrule
\end{tabular}
\end{center}
A Curva de Koch e ``mais complexa'' que uma linha ($D>1$), mas ainda ``cabe'' dentro
do plano sem preenche-lo ($D<2$).

\end{document}
